% ==============================================================================
% PAGE 13: CHAPTER-2 LITERATURE REVIEW (Numbered Page 5)
% ==============================================================================
\begin{center}
  {\fontsize{15}{18}\selectfont \bfseries \textcolor{red!80!black}{\underline{CHAPTER-2}}}\\[2mm]
  {\fontsize{16}{20}\selectfont \bfseries \textcolor{red!80!black}{\underline{Literature Review}}}\\[5mm]
\end{center}

\noindent
\begin{spacing}{1.25}
{\fontsize{10.2}{14.5}\selectfont
\textbf{Literature Survey} Solar energy has been harnessed for heating and power generation for several centuries, dating back to ancient civilizations that used mirrors and lenses to concentrate sunlight for practical applications. Early research on solar concentrators and parabolic reflectors in the 18th and 19th centuries laid the foundation for modern solar thermal systems. Initial designs were simple, often stationary reflectors used for heating water or cooking food, and lacked automation or efficiency optimization, but they demonstrated the potential of concentrated solar energy.\\[3.5mm]
In the latter half of the 20th century, advancements in materials, reflectors, and control systems enabled the development of more efficient parabolic troughs and dishes. Researchers focused on improving the concentration ratio, heat transfer efficiency, and receiver design. During this period, solar tracking mechanisms also began to emerge, initially as manual or semi-automatic systems, allowing solar concentrators to follow the sun's path and improve energy capture. These developments highlighted the importance of combining solar tracking with concentrators for optimal performance.\\[3.5mm]
In the present era, significant improvements in photovoltaic (PV) technology and automated solar tracking systems have made hybrid systems combining thermal and electrical energy feasible. Modern parabolic water heaters now integrate solar tracking, temperature sensors, and electric motors powered by solar panels, creating partially self-sustaining systems. Studies in recent years have shown that such hybrid systems can improve overall energy utilization by up to 30--50\% compared to fixed or non-tracked systems. Additionally, the use of readily available materials and cost-effective fabrication methods has made these technologies accessible for both domestic and industrial applications.\\[3.5mm]
Future trends in parabolic water heating and power generation focus on smart automation, AI-based tracking, and multi-functional hybrid systems. Researchers are exploring the integration of IoT devices for real-time monitoring, predictive control of solar tracking, and energy storage optimization. Advanced materials like selective coatings for receivers and high-reflectivity mirrors are also being developed to enhance heat absorption and reduce energy losses. These innovations aim to create more efficient, durable, and environmentally friendly solar systems capable of meeting increasing energy demands.
}
\end{spacing}
\clearpage
